% Bản 15 trang: 19 tài liệu. Bỏ BIPIA, Nasr 2025, Tramèr 2020, Dwork 2015, Recht 2019, Holm 1979. Giữ FIDES, design
% patterns và Liu et al. vì phản biện 3 coi việc cắt các tài liệu này là suy giảm.
% CapLite (arXiv:2611.17115, 12/11/2026): chỉ dùng mức tóm tắt đã đọc ngày 19/11 (facts.md, "Tài liệu mới được trích").
\section{Nghiên cứu liên quan}\label{sec:related}

\paragraph{Tấn công và phòng thủ dựa vào mô hình.} Liu và cộng sự~\cite{liu2024formalizing} hình thức hóa tấn công chèn
lệnh và phòng thủ, còn InjecAgent~\cite{zhan2024injecagent} và AgentDojo~\cite{debenedetti2024agentdojo} đo tấn công
trên tác tử dùng công cụ. Các bộ dữ liệu này đều bằng tiếng Anh, trong khi cơ chế an toàn của LLM yếu hơn ở ngôn ngữ ít
tài nguyên~\cite{deng2024multijail}, và tấn công thích nghi đã vượt qua nhiều phòng thủ đã công
bố~\cite{zhan2025adaptive}. Spotlighting~\cite{hines2024spotlighting} đánh dấu dữ liệu không đáng tin, lời nhắc kiểu
sandwich lặp lại yêu cầu sau mỗi quan sát, SecAlign~\cite{chen2025secalign} huấn luyện lại mô hình để tách chỉ thị khỏi
dữ liệu, Prompt Guard~2~\cite{meta2025promptguard2} sàng lọc đầu vào bằng một bộ phân loại, còn
MELON~\cite{zhu2025melon} so hành động của lượt chạy gốc với lượt chạy lại có che yêu cầu. Các cách này phụ thuộc vào
hành vi của mô hình hoặc bộ phân loại trên chính dữ liệu tấn công.

\paragraph{Cưỡng chế bên ngoài mô hình.} Progent~\cite{shi2025progent} áp chính sách đặc quyền lên lời gọi công cụ, còn
ClawGuard~\cite{zhao2026clawguard} nhờ mô hình suy ra luật truy cập từ yêu cầu của người dùng.
FIDES~\cite{costa2025fides} gắn nhãn tin cậy và bảo mật cho dữ liệu rồi lan truyền nhãn qua các lời gọi.
CaMeL~\cite{debenedetti2026camel} tách luồng điều khiển khỏi dữ liệu không đáng tin bằng một trình thông dịch riêng và
gắn capability cho từng giá trị. CapLite~\cite{xu2026caplite}, công bố cùng thời gian, suy chính sách từ schema công cụ,
cam kết các lời gọi ghi trước khi đọc kết quả công cụ và chỉ nhận lời gọi ghi khi các đối số quyết định đối tượng có
trong yêu cầu. Beurer-Kellner và cộng sự~\cite{beurerkellner2025patterns} hệ thống hóa các mẫu thiết kế kiểu này. Bản
thân việc dùng chính sách có cấu trúc hay chứng minh an toàn vì vậy chưa phải là đóng góp mới. \ours{} khác ở chỗ giữ
nguyên vòng lặp ReAct và chỉ kiểm từng lời gọi tại broker, thay vì viết lại tác tử thành trình thông dịch như CaMeL hay
cố định tập lời gọi ghi trước khi đọc dữ liệu như CapLite. Chính sách được trích tất định từ $u$ và xét cực tính trong
tiếng Việt có dấu, không dấu và tiếng Anh, trong khi ClawGuard để mô hình viết luật và Progent dựa vào luật viết tay. Độ
nhạy cảm được gán theo kiểu giá trị ngay tại broker, không lan truyền qua mọi phép tính như FIDES, nên rẻ hơn nhưng chỉ
bảo vệ giá trị trong một lớp biến đổi đã nêu (Mục~\ref{sec:formal}).
